\begin{resumo}
Este trabalho investiga a baixa cobertura de testes no AgroMart Web e seus impactos sobre a manutenibilidade e a evolução segura da aplicação.
O diagnóstico inicial, realizado com apoio do SonarQube, identificou cobertura de apenas 0,5\%, com 2186 linhas elegíveis e 14 linhas cobertas por testes,
indicando elevado risco de regressões durante alterações no código. Diante desse cenário, o objetivo do estudo é elevar a qualidade interna do projeto por meio
da ampliação progressiva de testes unitários em unidades fundamentais do sistema. A metodologia adotada é de natureza exploratória, estruturada em ciclos de diagnóstico,
priorização de unidades, implementação de testes, medição de métricas e consolidação de práticas em documentação de desenvolvimento.
Como resultado, o trabalho estabelece uma linha de base técnica e um fluxo de melhoria contínua orientado por evidências, integrando análise estática,
execução de testes e acompanhamento sistemático de indicadores de qualidade. Conclui-se que a estratégia proposta oferece um caminho viável para aumentar
a cobertura de testes de forma sustentável e reduzir riscos de manutenção no contexto de evolução colaborativa do AgroMart Web.

\vspace{\onelineskip}

\noindent
\textbf{Palavras-chave}: qualidade de software. cobertura de testes. testes unitários. AgroMart.
\end{resumo}
